
% !TEX program = pdflatex
% ============================================================
%  Final Exam (3 Calculation Questions)
%  Strictly following notation in:
%   - Julia K. Thomas E8722 Topic 4 (Endogenous Growth): Slides Topic 04 E8722.pdf
%   - Aubhik Khan (2019): production1b.pdf
%
%  Toggle solutions with \showAnstrue / \showAnsfalse
% ============================================================

\documentclass[12pt]{article}
\usepackage[margin=1in]{geometry}
\usepackage{amsmath,amssymb,mathtools}
\usepackage{enumitem}
\usepackage{comment}

% -----------------------
% Show / hide solutions
% -----------------------
\newif\ifshowAns
% \showAnstrue      % <-- instructor version
\showAnsfalse       % <-- student version

\ifshowAns
  \newenvironment{sol}{\par\medskip\noindent\textbf{Solution.}\ }{\par\medskip}
\else
  \excludecomment{sol}
\fi

\newcommand{\pts}[1]{\hfill{\small[#1 pts]}}
\newcommand{\E}{\mathbb{E}}

\begin{document}

\begin{center}
{\Large \textbf{Final Exam}}\\
PhD Macro\\
\medskip
\end{center}
\textbf{Instructions.} Show steps. Partial credit for correct intermediate steps.

% ============================================================
\section*{Question 1: human capital investment using goods \pts{35}}

In the ``human capital investment using goods'' model (Topic 4, Case 4), the planner chooses next-period stocks
\((K',X')\) and consumption \(C\) subject to
\[
C = F(K,X) + (1-\delta)K + (1-\delta_x)X - K' - X',
\]
with \(\delta\in(0,1)\) and \(\delta_x\in(0,1)\).
Assume Cobb--Douglas technology
\[
F(K,X)=K^{\alpha}X^{1-\alpha}, \qquad \alpha\in(0,1),
\]
and CRRA utility
\[
u(C)=\frac{C^{1-\sigma}-1}{1-\sigma},\qquad \sigma>0.
\]
Let \(k\equiv K/X\) and write \(F(K,X)=X f(k)\) where \(f(k)=k^{\alpha}\).

\begin{enumerate}[label=(\alph*)]
\item \textbf{BGP condition for \(k^\ast\).} \pts{14}\\
Using the slide result that an interior BGP with \(K'>0\) and \(X'>0\) implies
\[
f'(k_{t+1})-\delta = f(k_{t+1})-k_{t+1}f'(k_{t+1})-\delta_x,
\]
and imposing \(k_{t+1}=k_t=k^\ast\) on the BGP, show that \(k^\ast\) solves
\[
(1+k)f'(k)-f(k)=\delta-\delta_x.
\tag{3}
\]
For Cobb--Douglas \(f(k)=k^\alpha\), show that when \(\delta=\delta_x\) the unique solution is
\[
k^\ast=\frac{\alpha}{1-\alpha}.
\]
(If \(\delta\neq \delta_x\), \(k^\ast\) is obtained by solving (3) numerically, as done in the slides.)

\item \textbf{Growth rate of growing series.} \pts{10}\\
Using the slide formula for the growth rate of growing series,
\[
\gamma_Y = \cdots = \gamma_C =
\left[\beta\left(f'(k^\ast)+1-\delta\right)\right]^{1/\sigma},
\]
compute \(\gamma_Y\) for
\[
\alpha=0.30,\quad \delta=0.10,\quad \delta_x=0.15,\quad \beta=0.985,\quad \sigma=2,
\]
using the slide-reported solution \(k^\ast=0.42857\).

\item \textbf{Interpretation (one line).} \pts{11}\\
Holding \((\alpha,\beta,\sigma,\delta)\) fixed, what is the direction of change in \(k^\ast\) when \(\delta_x\) increases?
Briefly explain using equation (3).
\end{enumerate}

\begin{sol}
(a) Rearranging the slide equation gives:
\[
f'(k_{t+1})+k_{t+1}f'(k_{t+1})-f(k_{t+1})=\delta-\delta_x
\Rightarrow (1+k_{t+1})f'(k_{t+1})-f(k_{t+1})=\delta-\delta_x.
\]
On the BGP, \(k_{t+1}=k_t=k^\ast\), yielding (3).
For Cobb--Douglas, \(f'(k)=\alpha k^{\alpha-1}\). If \(\delta=\delta_x\), (3) becomes
\[
(1+k)\alpha k^{\alpha-1}-k^\alpha=0
\Rightarrow k^{\alpha-1}\big(\alpha-(1-\alpha)k\big)=0
\Rightarrow k^\ast=\frac{\alpha}{1-\alpha}.
\]

(b) With \(k^\ast=0.42857\), \(f'(k^\ast)=\alpha(k^\ast)^{\alpha-1}=0.30\cdot (0.42857)^{-0.70}\approx 0.543\).
Then \(f'(k^\ast)+1-\delta\approx 0.543+0.90=1.443\).
So
\[
\gamma_Y\approx [0.985\cdot 1.443]^{1/2}\approx [1.421]^{1/2}\approx 1.192.
\]

(c) Equation (3) has RHS \(\delta-\delta_x\). Increasing \(\delta_x\) lowers RHS.
In the Cobb--Douglas case, the solution shifts so that \(K/X\) rises: faster depreciation of \(X\) makes \(X\) relatively harder to accumulate,
so the balanced ratio features higher \(k^\ast=K/X\).
\end{sol}

% ============================================================
\section*{Question 2: Lucas (1988) human capital accumulation model \pts{30}}

In the Lucas (1988) model (Topic 4), specialize production to
\[
Y_t = K_t^{\alpha} M_t^{1-\alpha}, \qquad \alpha\in(0,1),
\]
where \(M_t=\phi_t X_t\), \(X_t\) is the stock of skills, and \(\phi_t\in(0,1)\) is the fraction of unit time endowment spent working.
The fraction \(1-\phi_t\) is time in education.
On a balanced growth path with constant \(\phi\), skills grow according to the slide equation:
\[
\gamma_X \equiv \frac{X_{t+1}}{X_t} = 1 + A(1-\phi),
\qquad A>0.
\]
Preferences are CRRA with parameter \(\sigma\) and discount factor \(\beta\).

\begin{enumerate}[label=(\alph*)]
\item \textbf{BGP growth rate.} \pts{12}\\
Using the slide result, state and use the key implication:
\[
\gamma_Y = \cdots = \gamma_X = \left[\beta(1+A)\right]^{1/\sigma}.
\]

\item \textbf{Solve for \(\phi\).} \pts{10}\\
Combine \(\gamma_X = 1 + A(1-\phi)\) with the BGP growth rate to show:
\[
\phi = \frac{1}{A}\left[1+A-\left(\beta(1+A)\right)^{1/\sigma}\right].
\]

\item \textbf{Numerical.} \pts{8}\\
Compute \(\phi\), \(\gamma_X\), and \(\gamma_Y\) when
\[
\beta=0.96,\quad \sigma=2,\quad A=0.20.
\]
\end{enumerate}

\begin{sol}
(a) Topic 4 derives: \(\gamma_Y=\cdots=\gamma_X=[\beta(1+A)]^{1/\sigma}\).

(b) Solve \(1+A(1-\phi)=[\beta(1+A)]^{1/\sigma}\):
\[
A(1-\phi)=[\beta(1+A)]^{1/\sigma}-1
\Rightarrow
\phi=\frac{1}{A}\left[1+A-[\beta(1+A)]^{1/\sigma}\right].
\]

(c) \([\beta(1+A)]^{1/\sigma}=[0.96\cdot 1.20]^{1/2}=[1.152]^{1/2}\approx 1.073.\)
Then \(\phi=\frac{1}{0.20}(1.20-1.073)=5(0.127)\approx 0.635.\)
So \(\gamma_X=1+A(1-\phi)=1+0.20(0.365)\approx 1.073\), and \(\gamma_Y=\gamma_X\approx 1.073.\)
\end{sol}

% ============================================================
\section*{Question 3: aggregation over heterogenous firms and \(A(\mu)\) \pts{35}}

There is a large number of heterogeneous firms. Production at each firm is
\[
y = z\varepsilon k^{\alpha}n^{\nu},
\tag{1}
\]
where \(\alpha\in(0,1)\), \(\nu\in(0,1)\), and \(\alpha+\nu<1\).
Capital is rented and labor is hired. Profits are
\[
\pi(z\varepsilon;r,w)\equiv \max_{k,n}\; z\varepsilon k^{\alpha}n^{\nu}-(r+\delta)k-wn,
\tag{2}
\]
where \(r\) is the real interest rate, \(w\) is the real wage, and \(\delta\in(0,1)\).
Optimal choices satisfy
\[
\alpha z\varepsilon k^{\alpha-1}n^{\nu}=r+\delta,
\qquad
\nu z\varepsilon k^{\alpha}n^{\nu-1}=w.
\]
Let \(\mu(\varepsilon)\) describe the measure of firms in production over \(\varepsilon\).

\begin{enumerate}[label=(\alph*)]
\item \textbf{Firm choices (equations (6)--(8)).} \pts{12}\\
Dividing (3) by (4) gives equation (6):
\[
\frac{\alpha}{\nu}\frac{n}{k}=\frac{r+\delta}{w}.
\tag{6}
\]
Using (6) in (3), show that capital demand is (equation (7) in the slides):
\[
k(\varepsilon)=
\left(
\frac{z\varepsilon\,\alpha^{1-\nu}\nu^{\nu}}{(r+\delta)^{1-\nu}w^{\nu}}
\right)^{\frac{1}{1-(\alpha+\nu)}}.
\tag{7}
\]
(You do not need to re-derive the labor demand formula (8).)

\item \textbf{Market clearing and the key moment.} \pts{12}\\
Using the slide market-clearing equations (10)--(11),
\[
K = \kappa_K(z,r+\delta,w)\int \varepsilon^{\frac{1}{1-(\alpha+\nu)}}\,\mu(d\varepsilon),
\tag{10}
\]
\[
N = \kappa_N(z,r+\delta,w)\int \varepsilon^{\frac{1}{1-(\alpha+\nu)}}\,\mu(d\varepsilon),
\tag{11}
\]
define
\[
\mathcal{M}(\mu)\equiv \int \varepsilon^{\frac{1}{1-(\alpha+\nu)}}\,\mu(d\varepsilon).
\]
Explain why, holding \((z,r+\delta,w)\) fixed, changes in \(\mu\) affect \(K\) and \(N\) only through \(\mathcal{M}(\mu)\).

\item \textbf{Distribution changes and aggregate productivity.} \pts{11}\\
Define an ``endogenous TFP'' index
\[
A(\mu)\equiv \left(\mathcal{M}(\mu)\right)^{\,1-(\alpha+\nu)}.
\]
Let \(\alpha=0.30\), \(\nu=0.50\) so \(1-(\alpha+\nu)=0.20\). Compare:
\[
\mu_1:\ \varepsilon=1\ \text{w.p. }1,
\qquad
\mu_2:\ \varepsilon=0.5\ \text{w.p. }0.5,\ \varepsilon=1.5\ \text{w.p. }0.5.
\]
Compute \(A(\mu_1)\) and \(A(\mu_2)\). Does greater dispersion raise or lower \(A(\mu)\) here?
\end{enumerate}

\begin{sol}
(a) This is exactly equation (7) in production1b.

(b) In (10)--(11), both \(K\) and \(N\) depend on \(\mu\) only through the same integral moment \(\mathcal{M}(\mu)\).
Hence shifts in the distribution that change \(\mathcal{M}(\mu)\) shift aggregate inputs (for fixed prices) and, equivalently,
shift measured productivity in the aggregate production function.

(c) Here \(\frac{1}{1-(\alpha+\nu)}=\frac{1}{0.20}=5\), so \(\mathcal{M}(\mu)=\E[\varepsilon^5]\) and \(A(\mu)=\mathcal{M}(\mu)^{0.20}\).
For \(\mu_1\): \(\mathcal{M}=1\Rightarrow A(\mu_1)=1\).
For \(\mu_2\): \(\mathcal{M}=0.5(0.5^5)+0.5(1.5^5)=0.5(0.03125+7.59375)=3.8125\).
So \(A(\mu_2)=3.8125^{0.20}\approx e^{0.20\ln 3.8125}\approx e^{0.268}\approx 1.307>1\).
Thus more dispersion raises \(A(\mu)\) in this example (convexity of \(\varepsilon^5\)).
\end{sol}

\end{document}
